\documentclass[10pt,twocolumn]{article}

\usepackage{fontspec}
\usepackage[margin=0.72in]{geometry}
\usepackage{microtype}
\usepackage{cite}
\usepackage{amsmath}
\usepackage{unicode-math}
\usepackage{graphicx}
\usepackage{booktabs}
\usepackage{multirow}
\usepackage{xcolor}
\usepackage{hyperref}
\usepackage{caption}

\setcounter{dbltopnumber}{4}
\renewcommand{\dbltopfraction}{0.88}
\renewcommand{\textfraction}{0.08}

\hypersetup{
    colorlinks=true,
    linkcolor=blue!60!black,
    citecolor=blue!60!black,
    urlcolor=blue!60!black
}

\title{\textbf{ABK-VLA: Task-ID-Free Continual Visuomotor Policy Learning via Autonomous B-Spline Knot Compression and Visual Prototype Routing}}

\author{
  Seungjoon Son \\
  Independent AI \& Robotics Researcher, Daejeon, Republic of Korea \\
  \texttt{https://github.com/5Color/abk-vla}
}

\date{November 2026}

\begin{document}

\maketitle

\begin{abstract}
Catastrophic forgetting remains a critical bottleneck for compact, real-time visuomotor policies deployed in sequential manipulation tasks. While large-scale foundation models mitigate forgetting through over-parameterization and cloud-scale data replay, lightweight closed-loop controllers suffer severe trajectory degradation when trained sequentially under strict memory and latency budgets. In this paper, we propose \textbf{ABK-VLA} (\textit{Autonomous B-spline Knot-based Visuomotor Policy}), evaluated on both a goal-image-conditioned 2D multi-stroke drawing testbed and a 4-task 3D Meta-World manipulation benchmark without experience replay buffers or inference-time \texttt{Task-ID} supervision. ABK-VLA compresses action chunks ($H=16$) into low-dimensional cubic B-spline control knots ($N_k=6$), allowing a frozen visual-spatial backbone ($4.59\text{M}$ shared parameters; $94.44\%$ of total $5.90\text{M}$ parameters frozen by $T_3$) to acquire new skills via compact 6-mode residual Winner-Take-All (WTA) slot heads ($0.33\text{M}$ trainable parameters per residual slot, $7.16\%$ of the backbone). At test time, a visual prototype router ($T=0.05$, zero task-routing heuristics) computes a $5{,}376$-dimensional hybrid Conv4-and-spatial target key and aggregates slots via \textit{Select-then-Blend} or \textit{Hard Argmax} routing. Across four sequential 2D tasks (apple $\to$ sprout $\to$ star $\to$ cat), ABK-VLA achieves \textbf{99.98\% average diagonal routing confidence}, reduces forgetting from $89.33\%$ (naive fine-tuning) and $14.13\%$ (replay) down to \textbf{1.66\%} (Select-then-Blend, $85.59\%$ $F_1$\,) and \textbf{0.00\%} (Hard Argmax, $84.39\%$ $F_1$\,), maintains $0.61\pm 0.74\%$ (Blend) and $0.00\pm 0.00\%$ (Argmax) forgetting across 3 seeds and reverse ordering, retains $100.0\%$ OOD routing up to $\pm 8\text{ px}$ shift / $\pm 15^\circ$ rotation / $\sigma=0.40$ noise, and transfers directly to 3D Meta-World manipulation with \textbf{91.7\% average success rate}, \textbf{0.0\% forgetting}, and \textbf{+11.1\% backward transfer}.
\end{abstract}

\section{Introduction}
Real-world visuomotor agents must incrementally acquire new manipulation skills without erasing previously mastered behaviors. In large vision-language-action (VLA) foundation models~\cite{kim2024openvla}, massive parameter capacity and large-batch continual pre-training often dilute catastrophic forgetting. However, closed-loop visuomotor control policies operating on resource-constrained hardware cannot rely on billion-parameter backbones or unbounded replay buffers. When compact sequence-to-trajectory controllers are updated sequentially on new tasks, gradient interference overwrites shared motor representations, causing catastrophic trajectory collapse.

Existing parameter-isolation methods prevent weight interference by freezing task-specific subnetworks~\cite{rusu2016,mallya2018,serra2018}, yet they typically suffer from two limitations: (1) high per-task memory overhead when duplicating dense action decoders, and (2) reliance on an external oracle \texttt{Task-ID} during inference to activate the correct module~\cite{farquhar2018}.

To address both barriers while establishing clean controlled baselines prior to language-backbone integration, we evaluate \textbf{ABK-VLA} as a modular vision-and-proprioception action-chunking architecture across 2D drawing and 3D manipulation:
\begin{itemize}
    \item \textbf{Compact B-Spline Parameter Isolation:} By representing $H=16$ continuous action chunks via $N_k=6$ cubic B-spline control knots, each new residual WTA slot requires $0.33\text{M}$ trainable parameters ($7.16\%$ of the $4.59\text{M}$ shared visual-state backbone), keeping $94.44\%$ of total model parameters frozen by $T_3$.
    \item \textbf{Task-ID-Free Visual Prototype Routing (Select-then-Blend):} During inference, ABK-VLA caches a $5{,}376$-dimensional hybrid Conv4-and-spatial target key per episode (\texttt{reset\_cache()}), computes temperature-scaled cosine similarity weights ($T=0.05$) without task-specific branching, and blends sequential residual slot outputs, achieving $99.98\%$ mean diagonal routing confidence across canonical target templates.
    \item \textbf{Horizon Alignment \& Post-Hoc Continuity:} Executing full $H=16$ chunks (\texttt{exec\_steps}=16) preserves pen-up jump landings across disconnected strokes, while post-hoc chunk-boundary $C^0/C^1$ smoothing yields low-jerk trajectories ($\mathrm{NMJ}=3.33$).
\end{itemize}

\section{Related Work}
\textbf{Continual Learning \& Parameter Isolation.} Continual learning methods traditionally mitigate catastrophic forgetting via parameter regularization (e.g., EWC~\cite{kirkpatrick2017}, Synaptic Intelligence~\cite{zenke2017}, and Learning without Forgetting~\cite{li2017}), experience replay buffers (e.g., ER~\cite{rolnick2019} and iCaRL~\cite{rebuffi2017}), or architectural parameter isolation (e.g., Progressive Neural Networks~\cite{rusu2016}, PackNet~\cite{mallya2018}, and Hard Attention to the Task~\cite{serra2018}); see Parisi et al.~\cite{parisi2019} and De Lange et al.~\cite{delange2022} for broader taxonomies. While parameter isolation eliminates weight overwriting, classical implementations require explicit test-time task identifiers~\cite{farquhar2018}.

\textbf{Task-ID-Free Prototype \& Prompt Routing.} Metric-based prototypical representations~\cite{snell2017} enable similarity-driven inference without fixed classifier heads. Recent rehearsal-free prompt-pool methods---such as L2P~\cite{wang2022l2p}, DualPrompt~\cite{wang2022dualprompt}, and CODA-Prompt~\cite{smith2023}---query small prompt components over frozen Vision Transformers without test-time \texttt{Task-ID}, while parameter-efficient adapters such as LoRA~\cite{hu2022lora} isolate task updates in low-rank subspaces. ABK-VLA brings key-query prototype routing to lightweight ConvNet closed-loop controllers by pairing $5{,}376$-d Conv4+spatial keys with logit-space residual B-spline slots.

\textbf{Continual Robot Learning, Action Chunking \& Splines.} Lifelong manipulation benchmarks including Meta-World~\cite{yu2020metaworld} and LIBERO~\cite{liu2023libero} highlight severe trajectory forgetting during sequential imitation learning, prompting adapter-based (TAIL~\cite{liu2024tail}) and skill-discovery (LOTUS~\cite{wan2024lotus}) continual policies. Temporally extended action chunking~\cite{zhao2023}, generative action diffusion~\cite{chi2023}, discretized multimodal policies (BeT~\cite{shafiullah2022}), and Winner-Take-All (WTA) multiple-choice predictors~\cite{guzman2012,rupprecht2017} mitigate compounding drift and trajectory multimodality. Building on classical spline and movement primitives---DMPs~\cite{ijspeert2013}, ProMPs~\cite{paraschos2013}, and clamped cubic B-splines~\cite{piegl1997}---ABK-VLA compresses $H=16$ action chunks into $N_k=6$ knots~\cite{piegl1997} so isolated multi-mode residual slots remain lightweight ($0.19\text{M}$--$0.33\text{M}$/slot).

\section{Methodology}

\subsection{Cubic B-Spline Trajectory Compression}
Given an action horizon $H=16$, rather than predicting $H$ independent $(x_t, y_t, p_t)$ waypoints, each slot $s \in \{0, \dots, K-1\}$ predicts $N_k=6$ spatial control knots $\mathbf{P}_{s,k} \in (0,1)^{N_k \times 2}$ (via sigmoid activation) across $M=6$ Winner-Take-All (WTA) modes $k \in \{0, \dots, M-1\}$~\cite{guzman2012,rupprecht2017}, together with step-wise pen-contact logits $\mathbf{u}_{s,k} \in \mathbb{R}^{H}$ and mode selection logits $\mathbf{g}_s \in \mathbb{R}^M$. Using a precomputed clamped cubic B-spline basis matrix $\mathbf{B} \in \mathbb{R}^{H \times N_k}$~\cite{piegl1997}, the continuous Cartesian trajectory $\hat{\mathbf{Y}} \in \mathbb{R}^{H \times 2}$ is reconstructed analytically via:
\begin{equation}
    \hat{\mathbf{Y}} = \mathbf{B}\mathbf{P},
\end{equation}
with a second-order control-point difference penalty $\mathcal{L}_{\text{smooth}} = \mathbb{E}[\|\mathbf{D}_2 \mathbf{P}\|_2^2]$ ($\lambda_{\text{sm}}=0.02$) suppressing intra-chunk loops. During closed-loop rollout, whenever both the previous executed step and the first predicted step are pen-down ($p_{t-1} > 0.5$ and $\hat{p}_0 > 0.5$), post-hoc boundary smoothing clamps $\hat{\mathbf{y}}_0 \leftarrow \mathbf{p}_{\text{prev}}$ ($C^0$) and blends $\hat{\mathbf{y}}_1 \leftarrow 0.7\,\hat{\mathbf{y}}_1 + 0.3\,(\mathbf{p}_{\text{prev}} + \mathbf{v}_{\text{prev}})$ whenever $\cos\,\alpha > 0$ ($C^1$ tangent continuity).

\subsection{Thin-Stroke Spatial Pointer \& Residual Cue}
To prevent thin 1\,px auxiliary strokes (e.g., stems, leaves, whiskers) from vanishing under spatial pooling, the policy derives two explicit residual cues from $\mathbf{M}_{\text{uncov}} = \max(0, \mathbf{I}_{\text{target}} - \mathbf{I}_{\text{canvas}})$:
\begin{enumerate}
    \item A $64$-dimensional residual map cue $\mathbf{c}_{\text{res}} = 4.0 \cdot \text{AvgPool}_{8\times 8}(\mathbf{M}_{\text{uncov}})$.
    \item A blended salience map for nearest-pixel localization:
    \begin{equation}
        \mathbf{M}_{\text{ptr}} = 0.7 \cdot \mathbf{M}_{\text{uncov}} + 0.3 \cdot \text{AvgPool}_{3\times 3}(\mathbf{M}_{\text{uncov}}),
    \end{equation}
    thresholded at $\mathbf{M}_{\text{ptr}} > 0.15$ (falling back to $\mathbf{M}_{\text{uncov}} > 0.08$ when empty) to locate the nearest uncovered pixel $(n_x, n_y)$, forming a $7$-dimensional geometric pointer:
    \begin{equation*}
        \mathbf{p}_{\text{geom}} = [c_x, c_y, r_{\text{rem}}, c_x - p_x, c_y - p_y, n_x - p_x, n_y - p_y].
    \end{equation*}
\end{enumerate}
The $4$-stage ConvNet visual encoder ($4{,}583{,}168$ params, $256$-d output) and $2$-layer MLP state encoder ($4{,}416$ params, $64$-d output) form the $4.59\text{M}$-parameter shared backbone, concatenated with $\mathbf{c}_{\text{res}}$ and $\mathbf{p}_{\text{geom}}$ into a $391$-dimensional conditioning vector $\mathbf{h}_{\text{cond}}$.

\subsection{Dynamic Residual Slot Expansion \& Select-then-Blend}
To resolve stroke-order multimodality and cross-task transfer, slot $0$ predicts base control knots directly while each expanded slot $s > 0$ (\texttt{ResidualWTASlotHead}, $328{,}367$ params) predicts logit-space residuals $(\Delta \mathbf{P}_s, \Delta \mathbf{u}_s)$ modulated by a learnable scalar gate $\alpha_s$ (\texttt{prior\_weight}, initialized at $1.0$):
\begin{equation}
    \tilde{\mathbf{P}}_s = \sigma\!\Big(\alpha_s \cdot \operatorname{logit}\big(\operatorname{clip}(\bar{\mathbf{P}}_{<s}, 0.05, 0.95)\big) + \Delta \mathbf{P}_s\Big), \label{eq:res_p}
\end{equation}
\begin{equation}
    \tilde{\mathbf{u}}_s = \alpha_s \cdot \operatorname{clip}(\bar{\mathbf{u}}_{<s}, -2.5, 2.5) + \Delta \mathbf{u}_s, \label{eq:res_u}
\end{equation}
where $(\bar{\mathbf{P}}_{<s}, \bar{\mathbf{u}}_{<s})$ is the softmax-weighted output of slots $0, \dots, s-1$. Each slot's mode selector \texttt{sel\_head} directly concatenates the $7$-d $\mathbf{p}_{\text{geom}}$ to its trunk features ($256 + 7 \rightarrow 128 \rightarrow M$).

\textbf{Hybrid 5,376-d Prototype Key \& Novelty Detection.} Given target image $\mathbf{I}_{\text{target}}$, \texttt{\_compute\_target\_key} extracts the $L_2$-normalized Conv4 spatial feature map ($4\times 4 \times 256 = 4{,}096$-d) weighted by $0.5$ and concatenates it with $L_2$-normalized multi-scale target silhouette average pooling ($32\times 32$ smoothed + $16\times 16 = 1{,}280$-d) weighted by $1.0$, followed by global $L_2$ normalization into $\mathbf{z} \in \mathbb{R}^{5{,}376}$. By construction, squared $L_2$ energy is split $80\% / 20\%$ between multi-scale spatial target pooling and Conv4 features. At task transitions during training, if $\max_{j < K} \mathbf{z}_\tau^\top \symbfit{\mu}_j < \gamma$ ($\gamma = 0.85$), a new residual slot $K$ is instantiated, $\symbfit{\mu}_K \leftarrow \mathbf{z}_\tau$ is registered, and the shared backbone ($4.59\text{M}$) plus all prior slots are permanently frozen (in \texttt{eval()} mode to lock BatchNorm statistics). After $T_0$, \texttt{refresh\_slot\_key} updates $\symbfit{\mu}_0$ prior to backbone freezing.

\textbf{Training Objective.} Active slot $s$ is trained without \texttt{Task-ID} using $M=6$ WTA supervision~\cite{rupprecht2017} over trajectory $L_1$, control-point $L_1$, smoothness, and transition-weighted pen BCE (\texttt{balanced\_pen\_transition\_loss}), combined with a $6.0\times$ inverse-frequency focal gating loss on air-flight/transition chunks and pen-consistent gating supervision:
\begin{equation}
    \mathcal{L}_D = 0.90\,\mathcal{L}_{\text{win}} + 0.05\,\mathcal{L}_{\text{all}} + 0.10\,\mathcal{L}_{\text{focal-gate}} + 0.10\,\mathcal{L}_{\text{pen-exp}}.
\end{equation}

\textbf{Select-then-Blend Inference.} At test time, without any \texttt{Task-ID}, routing weights $\mathbf{w} \in \Delta^{K-1}$ are computed once per episode via cached static keys (\texttt{reset\_cache()}, $T = 0.05$):
\begin{equation}
    w_s = \frac{\exp(\mathbf{z}^\top \symbfit{\mu}_s / T)}{\sum_{j=0}^{K-1} \exp(\mathbf{z}^\top \symbfit{\mu}_j / T)}.
\end{equation}
Rather than hard \texttt{argmax} routing over a single slot, ABK-VLA executes \textit{Select-then-Blend}: within each slot $s \in \{0, \dots, K-1\}$, \texttt{argmax} selects the best mode $k_s^* = \arg\max_k g_{s,k}$, residual composition chains sequentially across slots, and the final control knots and pen logits are blended across all slots as $\mathbf{P}^* = \operatorname{clip}(\sum_{s=0}^{K-1} w_s \tilde{\mathbf{P}}_{s,k_s^*}, 0, 1)$ and $\mathbf{u}^* = \sum_{s=0}^{K-1} w_s \tilde{\mathbf{u}}_{s,k_s^*}$.

\section{Experimental Setup}

\subsection{5-Model Continual Learning Protocol}
All experiments were conducted on Intel(R) Arc Graphics (\texttt{xpu}, FP32, PyTorch \texttt{2.15.0+xpu}, \texttt{seed=42}) across four sequential multi-stroke drawing tasks: $T_0(\text{apple}) \rightarrow T_1(\text{sprout}) \rightarrow T_2(\text{star}) \rightarrow T_3(\text{cat})$. Models~1--5 share identical episode-level train/val splits ($85\% / 15\%$), identical affine/elastic/covariate-shift augmentation schedules ($T_0$: $10\text{ epochs} \times 1\times$; $T_1$: $12\text{ epochs} \times 3\times$; $T_2$: $12\text{ epochs} \times 5\times$; $T_3$: $16\text{ epochs} \times 5\times$; batch size $64$; \texttt{WeightedRandomSampler} with $2.0\times$ secondary-stroke and $2.5\times$ jump-chunk weights on multi-stroke tasks), identical spatial pointers, post-hoc $C^0/C^1$ smoothing, and identical rollout parameters (\texttt{exec\_steps}=16, \texttt{max\_iters}=100, base \texttt{patience}=8, dynamically extended to \texttt{effective\_patience}=16 when $>5\%$ target ink remains, with stale accumulation throttled $1/3$ during air flights):
\begin{enumerate}
    \item \textbf{Model 1 (GRU + Naive FT):} Autoregressive $(x, y, p)$ GRUCell baseline trained sequentially at $\text{lr}=3\times 10^{-4}$ without isolation or replay.
    \item \textbf{Model 2 (ABK + Naive FT):} Single-mode B-spline MLP policy trained sequentially at $\text{lr}=3\times 10^{-4}$ with $\text{pen\_up\_weight}=2.0$ BCE.
    \item \textbf{Model 3 (ABK + Replay):} Model~2 augmented with a balanced replay buffer ($256$ samples/task) at $\text{lr}=3\times 10^{-4}$.
    \item \textbf{Model 4 (Oracle Slots / Model C):} Multi-slot 6-mode \texttt{WTASlotHead} ($0.30\text{M}$/slot) supplied with ground-truth \texttt{Task-ID} labels at test time ($\text{lr}=10^{-3}$ for $T_1$--$T_3$, cross-entropy gate loss).
    \item \textbf{Model 5 (Autonomous ABK-VLA / Model D, Ours):} Replay-free, \texttt{Task-ID}-free residual WTA policy ($0.33\text{M}$/slot, $\text{lr}=10^{-3}$ for $T_1$--$T_3$, $5\times 10^{-3}$ on $\alpha_s$) with focal and pen-consistent gating supervision.
\end{enumerate}
Regarding learning-rate selection, Models~1--3 update all $4.59\text{M}$ shared visual-backbone parameters across tasks and thus use $\text{lr}=3\times 10^{-4}$ for encoder stability, whereas Models~4--5 freeze the backbone after $T_0$ and optimize only compact $0.30\text{M}$--$0.33\text{M}$ MLP slot heads at $\text{lr}=10^{-3}$.

\subsection{Evaluation Metrics}
Let $R_{i,t}$ denote the pixel-level $F_1$ score ($\text{tol}=2$\,px) on task $i$ immediately after training on task $t$. We report standard Continual Learning Forgetting Rate ($\text{FR}$) and Backward Transfer ($\text{BWT}$):
\begin{equation}
    \text{FR} = \frac{1}{K-1} \sum_{i=0}^{K-2} \max\!\left(0,\; \max_{t \in \{i, \dots, K-2\}} R_{i,t} - R_{i, K-1}\right),
\end{equation}
\begin{equation}
    \text{BWT} = \frac{1}{K-1} \sum_{i=0}^{K-2} \big(R_{i, K-1} - R_{i,i}\big),
\end{equation}
alongside symmetric \textbf{Chamfer Distance (CD $\downarrow$)}, pen-down \textbf{Normalized Mean Jerk (NMJ $\downarrow$)}, connected \textbf{Component Recall (CR $\uparrow$)}, and per-episode \textbf{Inference Latency (ms $\downarrow$)}.

\section{Results and Diagnostic Analysis}

\begin{figure*}[t]
    \centering
    \includegraphics[width=\textwidth]{../outputs/autonomous_cl_rollout_grid.png}
    \caption{\textbf{Qualitative 5-model continual learning rollout comparison} after sequentially learning $\text{apple} \rightarrow \text{sprout} \rightarrow \text{star} \rightarrow \text{cat}$ (faint blue: target outline; black: executed ink). Models~1 and~2 suffer severe trajectory collapse on prior tasks, while Model~3 (Replay) fails on complex multi-part \texttt{Cat} geometry ($25.9\%$). Autonomous ABK-VLA (Model~5) preserves high structural fidelity across all four tasks without test-time \texttt{Task-ID} labels.}
    \label{fig:rollout_grid}
\end{figure*}

\begin{table*}[t]
\centering
\caption{5-Model Continual Learning Benchmark ($T_0$: apple $\to$ $T_1$: sprout $\to$ $T_2$: star $\to$ $T_3$: cat, \texttt{seed=42}).}
\label{tab:main_benchmark}
\resizebox{\textwidth}{!}{
\begin{tabular}{llccccc}
\toprule
\textbf{Axis} & \textbf{Metric} & \textbf{Model 1} & \textbf{Model 2} & \textbf{Model 3} & \textbf{Model 4 (Oracle)} & \textbf{Model 5 (Ours)} \\
 & & \texttt{GRU + Naive FT} & \texttt{ABK + Naive FT} & \texttt{ABK + Replay} & \texttt{ABK Slots (C)} & \texttt{Autonomous ABK (D)} \\
\midrule
\multirow{2}{*}{\textbf{Conditions}} & Inference \texttt{Task-ID} & None & None & None & Required (Oracle) & \textbf{None (Autonomous)} \\
 & Experience Replay & None & None & 256 / task & None & \textbf{None} \\
\midrule
\multirow{5}{*}{\textbf{Precision ($F_1 \uparrow$)}} & Apple ($T_0$) & 6.5\% & 4.7\% & 54.4\% & \textbf{99.7\%} & 88.5\% \\
 & Sprout ($T_1$) & 11.7\% & 5.4\% & 88.8\% & \textbf{92.2\%} & 81.1\% \\
 & Star ($T_2$) & 28.7\% & 8.3\% & 82.4\% & 87.1\% & \textbf{96.4\%} \\
 & Cat ($T_3$) & 1.6\% & 68.0\% & 25.9\% & \textbf{92.9\%} & 76.5\% \\
 & \textbf{Avg Final $F_1 \uparrow$} & 12.12\% & 21.61\% & 62.87\% & \textbf{92.96\%} & \textbf{85.59\%} \\
\midrule
\multirow{2}{*}{\textbf{Continual Learning}} & \textbf{Forgetting Rate $\downarrow$} & 46.44\% & 89.33\% & 14.13\% & \textbf{0.00\%} & \textbf{1.66\%} \\
 & Backward Transfer (BWT) & -46.44\% & -89.33\% & +1.95\% & 0.00\% & \textbf{-1.66\%} \\
\midrule
\multirow{3}{*}{\textbf{Geometry \& Control}} & \textbf{Chamfer Distance $\downarrow$} & 0.1043 & 0.0984 & 0.0327 & \textbf{0.0098} & \textbf{0.0192} \\
 & \textbf{Normalized Mean Jerk $\downarrow$} & 5.05 & 5.91 & 5.92 & \textbf{3.09} & \textbf{3.33} \\
 & \textbf{Component Recall $\uparrow$} & 0.0\% & 6.3\% & 50.0\% & \textbf{87.5\%} & \textbf{62.5\%} \\
\midrule
\multirow{2}{*}{\textbf{Efficiency}} & Episode Latency (ms) $\downarrow$ & 601.6 ms & 108.6 ms & 564.9 ms & 346.6 ms & \textbf{464.3 ms} \\
 & $T_3$ Trainable Params & 5.08M & 4.72M & 4.72M & 0.30M & \textbf{0.33M (94.44\% Frozen)} \\
\bottomrule
\end{tabular}
}
\end{table*}

\begin{figure*}[t]
    \centering
    \includegraphics[width=0.98\textwidth]{../outputs/autonomous_cl_heatmap.png}
    \caption{\textbf{Sequential $F_1$ retention heatmaps} across the 4-task curriculum for all five models.}
    \label{fig:heatmap}
\end{figure*}

\subsection{Quantitative Comparison}
As shown in Table~\ref{tab:main_benchmark} and Figures~\ref{fig:rollout_grid}--\ref{fig:heatmap}, sequential fine-tuning without isolation suffers severe forgetting: Model~1 (\texttt{GRU + Naive FT}) drops to $12.12\%$ average $F_1$ ($46.44\%$ forgetting, $0.0\%$ component recall), and Model~2 (\texttt{ABK + Naive FT}) exhibits $89.33\%$ forgetting. Experience replay (Model~3) mitigates forgetting ($14.13\%$), but its single-mode predictor trained at $\text{lr}=3\times 10^{-4}$ degrades to $25.9\%$ on multi-part \texttt{Cat} ($62.87\%$ average $F_1$).

In contrast, \textbf{Autonomous ABK-VLA (Model~5)} achieves an average final $F_1$ of \textbf{85.59\%}, recovering \textbf{92.1\%} of the external-\texttt{Task-ID} Model~4 ($92.96\%$) with a forgetting rate of \textbf{1.66\%} ($54.0\times$ lower than Model~2). Geometrically, Model~5 lowers Chamfer Distance to \textbf{0.0192} ($5.4\times$ better than GRU) with smooth trajectories ($\text{NMJ}=3.33$). Note that Model~5 outperforms Model~4 on \texttt{Star} ($96.4\%$ vs.\ $87.1\%$) due to its pointer-conditioned mode selector, focal gating, and residual head ($0.33\text{M}$ vs.\ $0.30\text{M}$). Because Select-then-Blend averages all slots with small residual tail weights ($<0.04\%$) rather than hard \texttt{argmax} gating, adding later slots slightly perturbs frozen earlier rollouts (\texttt{Apple}: $92.48\% \rightarrow 88.46\%$; \texttt{Star}: $97.30\% \rightarrow 96.36\%$), explaining the minor $1.66\%$ forgetting.

\subsection{Autonomous Visual Routing \& Latency Analysis}
\begin{figure}[t]
    \centering
    \includegraphics[width=0.88\columnwidth]{../outputs/autonomous_routing_matrix.png}
    \caption{\textbf{Task-ID-free softmax routing weights} ($T=0.05$) on canonical target templates.}
    \label{fig:routing}
\end{figure}
Evaluating the temperature-scaled softmax router ($T=0.05$, Figure~\ref{fig:routing}) on the four canonical task target templates yields diagonal routing weights of \texttt{Apple} $\rightarrow$ Slot~0 ($99.99\%$), \texttt{Sprout} $\rightarrow$ Slot~1 ($100.00\%$), \texttt{Star} $\rightarrow$ Slot~2 ($99.97\%$), and \texttt{Cat} $\rightarrow$ Slot~3 ($99.96\%$), giving \textbf{99.98\% mean diagonal confidence} ($4/4$ slot match). At task transitions, novelty detection ($\gamma=0.85$) cleanly triggered expansion with maximum prior similarities of $0.351$ (\texttt{sprout}), $0.491$ (\texttt{star}), and $0.593$ (\texttt{cat}). With static target key caching (\texttt{reset\_cache()}), Model~5 averages \textbf{464.3 ms} per episode ($7.02\text{ ms}$/chunk)---$1.30\times$ faster than GRU ($601.6\text{ ms}$) and $1.22\times$ faster than Replay ($564.9\text{ ms}$), while slower than single-slot Model~4 ($346.6\text{ ms}$) and early-stalling Model~2 ($108.6\text{ ms}$) because Select-then-Blend forwards all $K=4$ slots.

\subsection{Horizon Alignment \& Stall Patience Trade-Off}
Aligning $\text{exec\_steps}=16$ to the full $H=16$ spline horizon preserves pen-up jump landings without truncating airborne chunks prior to touchdown. Under the rollout stall timer (base $\text{patience}=8$, dynamically extended to $16$ when $>5\%$ target ink remains), long air-jumps toward thin auxiliary strokes (e.g., apple leaf, sprout stem, cat whiskers) still occasionally hit stall termination prior to landing ($62.5\%$ component recall vs.\ $87.5\%$ Oracle).

\subsection{Follow-Up Validation: Routing Stress Tests, Multi-Seed Order, Hard Argmax, and 3D Meta-World}

To rigorously stress-test routing invariance, stochastic seed variance, tail-weight forgetting, and 3D transfer, we conducted four additional controlled benchmarks.

\paragraph{Experiment 1: Perturbation Routing Robustness and Key Weight Ablation.}
We evaluated top-1 routing accuracy across 80 held-out perturbed queries per severity level (20 draws $\times$ 4 tasks) under five $(w_{\text{conv}}, w_{\text{sp}})$ key energy mixtures (\texttt{Spatial-Only} 0/100, \texttt{ABK Default} 20/80, \texttt{Balanced} 50/50, \texttt{Conv4-Heavy} 80/20, and \texttt{Conv4-Only} 100/0) and measured closed-loop rollout F1 under representative out-of-distribution (OOD) target corruptions (Fig.~\ref{fig:exp1_perturb}). The default 20/80 hybrid key maintains \textbf{100.0\%} routing accuracy up to $\pm 8$\,px translation, $\pm 15^\circ$ rotation, $\pm 15\%$ scale, $\sigma = 0.40$ Gaussian noise, and 20\% block occlusion (\textbf{92.5\%} at 40\% occlusion). Weighting Conv4 semantics (\texttt{Balanced} 50/50) lifts $\pm 12$\,px shift accuracy from 83.8\% to \textbf{90.0\%} (and \texttt{Conv4-Only} reaches 100.0\% at $\pm 12$\,px), whereas spatial pooling anchors invariance under noise (100.0\% vs.\ 50.0\% at $\sigma=0.40$) and occlusion (92.5\% vs.\ 70.0\% at 40\%). Under closed-loop OOD rollouts, Blend/Argmax retain 85.59\%/84.39\% clean, 65.04\%/56.97\% at 4\,px shift, 72.78\%/72.59\% at $8^\circ$ rotation, 29.55\%/22.49\% at $\sigma=0.15$ noise, and 42.48\%/42.85\% at 20\% occlusion. Notably, this reveals a clean decoupling between prototype routing and policy decoding: while the router still selects the exact slot with 100.0\% accuracy under 4\,px shifts and $\sigma=0.15$ noise, the frozen Conv4 visual encoder (which receives the target channel without sensor-noise augmentation during $T_0$ training) under-tracks fine pixel contours once $\mathbf{I}_{\text{target}}$ is heavily corrupted.

\begin{figure*}[t]
\centering
\includegraphics[width=0.90\textwidth]{../outputs/exp1_routing_perturbation.png}
\caption{\textbf{Experiment 1: Target Perturbation Routing Curves and Key Weight Ablation.} Hybrid Conv4+Spatial pooling balances spatial shift tolerance with noise, rotation, and occlusion invariance.}
\label{fig:exp1_perturb}
\end{figure*}

\paragraph{Experiment 2: Multi-Seed ($n=3$) and Task-Order Robustness.}
We trained and evaluated Autonomous ABK-VLA across three seeds ($\{42, 43, 44\}$) on both \textbf{Forward} ($\text{apple}\rightarrow\text{sprout}\rightarrow\text{star}\rightarrow\text{cat}$) and \textbf{Reverse} ($\text{cat}\rightarrow\text{star}\rightarrow\text{sprout}\rightarrow\text{apple}$) curricula (Table~\ref{tab:exp2_seed}, Fig.~\ref{fig:exp2_seed}). Across all six runs (24 task-slot transitions), autonomous novelty detection expanded with \textbf{99.98\%} routing diagonal confidence. Forward Select-then-Blend reaches $85.15 \pm 1.06\%$ F1 ($0.61 \pm 0.74\%$ FR), while Hard Argmax guarantees strictly $0.00 \pm 0.00\%$ FR at $84.58 \pm 1.14\%$ F1. Under Reverse ordering (starting from the complex multi-stroke \texttt{cat} task), forgetting remains near zero ($0.06 \pm 0.06\%$ Blend, $0.00 \pm 0.00\%$ Argmax).

\begin{table*}[t]
\centering
\caption{\textbf{Experiment 2: 3-Seed ($\{42,43,44\}$) Forward and Reverse Task-Order Summary (Mean $\pm$ Std).}}
\label{tab:exp2_seed}
\small
\setlength{\tabcolsep}{5pt}
\renewcommand{\arraystretch}{1.15}
\begin{tabular}{llcccc}
\toprule
\textbf{Task Order} & \textbf{Routing Mode} & \textbf{Avg $F_1$ $\uparrow$} & \textbf{FR $\downarrow$} & \textbf{CD $\downarrow$} & \textbf{CR $\uparrow$} \\
\midrule
Forward & Select-then-Blend ($T{=}0.05$) & \textbf{85.15 $\pm$ 1.06\%} & 0.61 $\pm$ 0.74\% & \textbf{0.0166 $\pm$ 0.0032} & \textbf{72.92 $\pm$ 7.80\%} \\
Forward & Hard Argmax & 84.58 $\pm$ 1.14\% & \textbf{0.00 $\pm$ 0.00\%} & 0.0168 $\pm$ 0.0036 & 70.83 $\pm$ 10.62\% \\
Reverse & Select-then-Blend ($T{=}0.05$) & 79.50 $\pm$ 4.97\% & 0.06 $\pm$ 0.06\% & 0.0280 $\pm$ 0.0112 & 77.08 $\pm$ 2.95\% \\
Reverse & Hard Argmax & 78.16 $\pm$ 5.53\% & \textbf{0.00 $\pm$ 0.00\%} & 0.0285 $\pm$ 0.0114 & 72.92 $\pm$ 7.80\% \\
\bottomrule
\end{tabular}
\end{table*}

\begin{figure*}[t]
\centering
\includegraphics[width=0.86\textwidth]{../outputs/exp2_multiseed_order.png}
\caption{\textbf{Experiment 2: Multi-Seed Forward vs.\ Reverse Task Order.} Error bars show std across seeds $\{42,43,44\}$; Hard Argmax eliminates forgetting in every seed and ordering.}
\label{fig:exp2_seed}
\end{figure*}

\paragraph{Experiment 3: Select-then-Blend Temperature Sweep vs.\ Hard Argmax.}
To verify the exact source of the remaining 1.66\% forgetting in the seed-42 benchmark, we re-evaluated the exact final checkpoint across softmax temperatures $T \in \{0.01, 0.05, 0.10, 0.25, 0.50\}$ and \texttt{hard\_argmax} (Table~\ref{tab:exp3_blend}, Fig.~\ref{fig:exp3_blend}). Switching strictly to \texttt{hard\_argmax} drops forgetting from \textbf{1.66\% to 0.00\%} and reduces standalone episode latency from \textbf{387.8\,ms to 357.5\,ms}, proving that all remaining drift arose exclusively from tiny off-diagonal residual weights ($<0.04\%$) rather than parameter corruption. Conversely, tuned soft blending ($T=0.05$) yields a $+1.20$\,pp net positive transfer boost on final F1 (85.59\% vs.\ 84.39\%, reaching 76.46\% vs.\ 66.69\% on \texttt{cat}), while over-smoothing ($T \ge 0.25$) collapses performance (FR $>78\%$). Notably, even under \texttt{Hard Argmax} (selecting a single slot $s^*$), Eqs.~(\ref{eq:res_p})--(\ref{eq:res_u}) still chain prior slots sequentially across $0 \dots s^*$ via $(\bar{\mathbf{P}}_{<s}, \bar{\mathbf{u}}_{<s})$, explaining why \texttt{Hard Argmax} ($357.5\text{ ms}$) reduces latency by only $7.8\%$ relative to \texttt{Blend (T=0.05)} ($387.8\text{ ms}$) and grows as $O(K)$ in compute. Regarding table cross-comparison, Table~\ref{tab:main_benchmark} reports $464.3\text{ ms}$ timed inline during the full 5-model sequential training-and-evaluation runner (\texttt{benchmark\_autonomous\_cl.py}), whereas Table~\ref{tab:exp3_blend} profiles the saved \texttt{autonomous\_model\_d\_final.pth} checkpoint in standalone isolation (\texttt{eval\_argmax\_vs\_blend.py}), giving apples-to-apples timing between \texttt{Blend (T=0.05)} and \texttt{Hard Argmax}; shorter times at $T \ge 0.25$ ($167.4\text{ ms}$ and $242.8\text{ ms}$) reflect premature rollout stall termination under trajectory collapse rather than compute speedup.

\begin{table*}[t]
\centering
\caption{\textbf{Experiment 3: Select-then-Blend vs.\ Hard Argmax on \texttt{autonomous\_model\_d\_final.pth}.}}
\label{tab:exp3_blend}
\small
\setlength{\tabcolsep}{7pt}
\renewcommand{\arraystretch}{1.15}
\begin{tabular}{lccccc}
\toprule
\textbf{Mode / Temperature} & \textbf{Avg $F_1$ $\uparrow$} & \textbf{FR $\downarrow$} & \textbf{CD $\downarrow$} & \textbf{CR $\uparrow$} & \textbf{Latency $\downarrow$} \\
\midrule
Blend ($T=0.01$) & 78.16\% & 6.25\% & 0.0232 & 56.2\% & 754.1\,ms \\
Blend ($T=0.05$, Default) & \textbf{85.59\%} & 1.66\% & \textbf{0.0192} & 62.5\% & 387.8\,ms \\
Blend ($T=0.10$) & 80.95\% & 3.66\% & 0.0226 & \textbf{68.8\%} & 365.4\,ms \\
Blend ($T=0.25$) & 10.70\% & 78.23\% & 0.1188 & 0.0\% & 167.4\,ms \\
Blend ($T=0.50$) & 5.21\% & 84.33\% & 0.1288 & 0.0\% & 242.8\,ms \\
\textbf{Hard Argmax} & 84.39\% & \textbf{0.00\%} & 0.0201 & 56.2\% & \textbf{357.5\,ms} \\
\bottomrule
\end{tabular}
\end{table*}

\begin{figure*}[t]
\centering
\includegraphics[width=0.84\textwidth]{../outputs/exp3_argmax_vs_blend.png}
\caption{\textbf{Experiment 3: Softmax Temperature Sweep vs.\ Hard Argmax.} Hard Argmax strictly enforces 0.00\% FR, while $T=0.05$ Blend trades 1.66\% drift for higher forward transfer.}
\label{fig:exp3_blend}
\end{figure*}

\paragraph{Experiment 4: 3D Meta-World Continual Manipulation Generalization.}
We ported the identical architecture---cubic B-spline chunking ($H=16 \rightarrow N_k=6$ across 4-DoF EEF+gripper actions, exec horizon 8), frozen Conv4+state trunk (91.45\% frozen at Task 3), logit-space Residual WTA slots, and autonomous prototype routing---to a 4-task Meta-World MuJoCo benchmark~\cite{yu2020metaworld} (\texttt{button-press-v3} $\rightarrow$ \texttt{door-open-v3} $\rightarrow$ \texttt{window-open-v3} $\rightarrow$ \texttt{faucet-open-v3}; 30 demos/task on seeds 0--34, 15 unseen initializations/task on seeds 35--49). Because static tabletop camera frames (\texttt{corner}, $128\times 128$ RGB) share $\sim$80\% background pixels, raw spatial keys saturate above $0.99$ across distinct tasks; subtracting a $15\times 15$ local average (DoG-like high-pass residual) before spatial pooling and storing $M_{\text{ex}}=8$ initial-frame exemplars per slot ($7,936$-d key, $T=0.003$, threshold $0.996$) sharply separates task transitions ($\max$ prior similarity rising cleanly from $0.8143 \rightarrow 0.9592 \rightarrow 0.9792$). Across all 60 held-out evaluation episodes (Table~\ref{tab:exp4_mw}, Fig.~\ref{fig:exp4_mw}), Autonomous 3D ABK-VLA achieves \textbf{100.0\%} autonomous routing accuracy without test-time \texttt{Task-ID} labels (99.71\% diagonal confidence), \textbf{0.0\%} forgetting, and \textbf{91.7\%} final average success rate (\textbf{+11.1\% BWT}), outperforming Naive FT (43.3\% SR, 75.6\% FR), 256-chunk Experience Replay (88.3\% SR, 13.3\% FR), and even Oracle Task-ID Slots (90.0\% SR, 0.0\% FR), while Hard Argmax attains 86.7\% SR with 0.0\% FR. Notably, in 3D manipulation, \texttt{Blend}'s $+11.1\%$ positive backward transfer stems from \texttt{window-open-v3} improving from $40.0\%$ immediately after $T_2$ to $73.3\%$ after acquiring lateral turning kinematics on $T_3$ (\texttt{faucet-open-v3}). While temperature scaling ($T=0.003$) expands the $T_3$ logit gap $(0.996 - 0.9792)/0.003 = 5.60$ ($>270\times$ softmax separation), the raw cosine margin narrows to $0.0168$ due to shared tabletop structure; scaling beyond $K=4$ tasks or discriminating tasks with identical $t=0$ reset geometry will therefore require foreground segmentation or goal/language conditioning.

\begin{table*}[t]
\centering
\caption{\textbf{Experiment 4: Meta-World 3D Continual Manipulation Benchmark (15 Held-Out Seeds/Task).}}
\label{tab:exp4_mw}
\small
\setlength{\tabcolsep}{7pt}
\renewcommand{\arraystretch}{1.15}
\begin{tabular}{lcccc}
\toprule
\textbf{Method} & \textbf{Test ID?} & \textbf{Avg SR $\uparrow$} & \textbf{FR $\downarrow$} & \textbf{BWT $\uparrow$} \\
\midrule
Naive FT (1-Head B-Spline) & No & 43.3\% & 75.6\% & -75.6\% \\
Experience Replay (256/task) & No & 88.3\% & 13.3\% & -13.3\% \\
Oracle Task-ID Slots & Yes & 90.0\% & \textbf{0.0\%} & +0.0\% \\
Autonomous 3D ABK (Argmax) & No & 86.7\% & \textbf{0.0\%} & +0.0\% \\
\textbf{Autonomous 3D ABK (Blend)} & \textbf{No} & \textbf{91.7\%} & \textbf{0.0\%} & \textbf{+11.1\%} \\
\bottomrule
\end{tabular}
\end{table*}

\begin{figure*}[t]
\centering
\includegraphics[width=0.78\textwidth]{../outputs/exp4_metaworld_cl.png}
\caption{\textbf{Experiment 4: Meta-World 3D Continual Manipulation.} Without test-time \texttt{Task-ID} or replay, Autonomous 3D ABK-VLA achieves 100.0\% autonomous routing accuracy, 0.0\% FR, and 91.7\% Avg SR.}
\label{fig:exp4_mw}
\end{figure*}

\section{Scope, Limitations, and Future Work}
We explicitly state four scope boundaries of this study:
\begin{enumerate}
    \item \textbf{Vision-and-Proprioception Simulation Scope:} ABK-VLA here denotes a compact vision-and-proprioception action-chunking architecture validated on a 2D drawing testbed and a 3D Meta-World (\texttt{MuJoCo}) continual manipulation benchmark~\cite{yu2020metaworld}; the current implementation contains no language encoder/VLM backbone~\cite{kim2024openvla} or physical robot hardware experiments.
    \item \textbf{Training Task Boundaries:} While inference is \texttt{Task-ID}-free, sequential training invokes \texttt{register\_or\_expand\_slot} at discrete task loader boundaries; autonomous online similarity-drop slot creation ($\max_j \mathbf{z}^\top \symbfit{\mu}_j < \gamma$) during unlabeled streams remains future work.
    \item \textbf{OOD Robustness, Narrow 3D Margin \& $t=0$ Identifiability:} In 2D, the hybrid key sustains 100.0\% routing up to $\pm 8$\,px shift, $\pm 15^\circ$ rotation, $\pm 15\%$ scale, $\sigma=0.40$ noise, and 20\% occlusion, while closed-loop OOD rollout $F_1$ degrades under unaugmented extreme corruptions ($65.04\%$ at 4\,px shift, $29.55\%$ at $\sigma=0.15$ noise). In 3D Meta-World, static background dominance required $15\times 15$ DoG high-pass subtraction and 8 initial-frame exemplars per slot, yielding a raw $T_3$ cosine margin of $0.0168$ ($0.9792$ vs.\ threshold $0.996$). Moreover, initial-frame ($t=0$) key caching assumes distinct initial scene layouts and cannot by itself disambiguate tasks sharing identical reset observations without goal-image or language conditioning.
    \item \textbf{$O(K)$ Memory/Compute Chaining \& Blend vs.\ Argmax Trade-Off:} Both slot memory ($0.33\text{M}$/slot in 2D, $5.90\text{M}$ at $T_3$; $0.19\text{M}$/slot in 3D, $2.19\text{M}$ at $T_3$) and sequential residual evaluation via Eqs.~(\ref{eq:res_p})--(\ref{eq:res_u}) scale linearly as $O(K)$, even under \texttt{hard\_argmax}. Across Experiments 2--3, \texttt{hard\_argmax} strictly eliminates forgetting ($0.00\%$ FR across all 3 seeds and both task orders) with a minor $-1.20$\,pp transfer gap relative to soft $T=0.05$ Select-then-Blend ($84.39\%$ vs.\ $85.59\%$ $F_1$); pruning low-$\alpha_s$ chains for $O(1)$ test-time slot execution is a promising extension.
\end{enumerate}

\section{Conclusion}
ABK-VLA demonstrates that pairing cubic B-spline trajectory compression with logit-space residual WTA slots and visual prototype routing enables compact visuomotor policies to eliminate catastrophic forgetting without replay buffers or test-time \texttt{Task-ID} supervision---achieving $85.59\%$ $F_1$ ($1.66\%$ FR) / $84.39\%$ $F_1$ ($0.00\%$ FR via \texttt{hard\_argmax}) in 2D continual drawing across 3 seeds and reverse task orders, and transferring directly to 3D Meta-World manipulation at $91.7\%$ average success rate, $0.0\%$ forgetting, $+11.1\%$ backward transfer, and $100.0\%$ autonomous routing accuracy without test-time \texttt{Task-ID}.

{\small
\begin{thebibliography}{29}
\bibitem{kim2024openvla} M. J. Kim et al., ``OpenVLA: An open-source vision-language-action model,'' \textit{CoRL} / \textit{arXiv:2406.09246}, 2024.
\bibitem{rusu2016} A. A. Rusu et al., ``Progressive neural networks,'' \textit{arXiv:1606.04671}, 2016.
\bibitem{mallya2018} A. Mallya and S. Lazebnik, ``PackNet: Adding multiple tasks to a single network by iterative pruning,'' \textit{CVPR}, 2018.
\bibitem{serra2018} J. Serr\`{a}, D. Sur\'{i}s, M. Miron, and A. Karatzoglou, ``Overcoming catastrophic forgetting with hard attention to the task,'' \textit{ICML}, 2018.
\bibitem{farquhar2018} S. Farquhar and Y. Gal, ``Towards robust evaluations of continual learning,'' \textit{arXiv:1805.09733}, 2018.
\bibitem{kirkpatrick2017} J. Kirkpatrick et al., ``Overcoming catastrophic forgetting in neural networks,'' \textit{PNAS}, vol. 114, no. 13, pp. 3521--3526, 2017.
\bibitem{zenke2017} F. Zenke, B. Poole, and S. Ganguli, ``Continual learning through synaptic intelligence,'' \textit{ICML}, 2017.
\bibitem{li2017} Z. Li and D. Hoiem, ``Learning without forgetting,'' \textit{IEEE TPAMI}, 2017.
\bibitem{rolnick2019} D. Rolnick, A. Ahuja, J. Schwarz, T. Lillicrap, and G. Wayne, ``Experience replay for continual learning,'' \textit{NeurIPS}, 2019.
\bibitem{rebuffi2017} S.-A. Rebuffi, A. Kolesnikov, G. Sperl, and C. H. Lampert, ``iCaRL: Incremental classifier and representation learning,'' \textit{CVPR}, 2017.
\bibitem{parisi2019} G. I. Parisi, R. Kemker, J. L. Part, C. Kanan, and S. Wermter, ``Continual lifelong learning with neural networks: A review,'' \textit{Neural Networks}, vol. 113, pp. 54--71, 2019.
\bibitem{delange2022} M. De Lange et al., ``A continual learning survey: Defying forgetting in classification tasks,'' \textit{IEEE TPAMI}, vol. 44, no. 7, pp. 3366--3385, 2022.
\bibitem{snell2017} J. Snell, K. Swersky, and R. S. Zemel, ``Prototypical networks for few-shot learning,'' \textit{NIPS}, 2017.
\bibitem{wang2022l2p} Z. Wang et al., ``Learning to prompt for continual learning,'' \textit{CVPR}, 2022.
\bibitem{wang2022dualprompt} Z. Wang et al., ``DualPrompt: Complementary prompting for rehearsal-free continual learning,'' \textit{ECCV}, 2022.
\bibitem{smith2023} J. S. Smith et al., ``CODA-Prompt: COntinual decomposed attention-based prompting for rehearsal-free continual learning,'' \textit{CVPR}, 2023.
\bibitem{hu2022lora} E. J. Hu et al., ``LoRA: Low-rank adaptation of large language models,'' \textit{ICLR}, 2022.
\bibitem{yu2020metaworld} T. Yu et al., ``Meta-World: A benchmark and evaluation for multi-task and meta reinforcement learning,'' \textit{CoRL}, 2019 / \textit{PMLR}, 2020.
\bibitem{liu2023libero} B. Liu et al., ``LIBERO: Benchmarking knowledge transfer for lifelong robot learning,'' \textit{NeurIPS Datasets and Benchmarks}, 2023.
\bibitem{liu2024tail} Z. Liu et al., ``TAIL: Task-specific adapters for imitation learning with large pretrained models,'' \textit{ICLR}, 2024.
\bibitem{wan2024lotus} W. Wan, Y. Zhu, R. Shah, and Y. Zhu, ``LOTUS: Continual imitation learning for robot manipulation through unsupervised skill discovery,'' \textit{ICRA}, 2024.
\bibitem{zhao2023} T. Z. Zhao, V. Kumar, S. Levine, and C. Finn, ``Learning fine-grained bimanual manipulation with low-cost hardware,'' \textit{RSS}, 2023.
\bibitem{chi2023} C. Chi et al., ``Diffusion Policy: Visuomotor policy learning via action diffusion,'' \textit{RSS}, 2023.
\bibitem{shafiullah2022} N. M. M. Shafiullah, Z. Cui, A. Altanzaya, and L. Pinto, ``Behavior Transformers: Cloning $k$ modes with one stone,'' \textit{NeurIPS}, 2022.
\bibitem{guzman2012} A. Guzm\'{a}n-Rivera, D. Batra, and P. Kohli, ``Multiple choice learning: Learning to produce multiple structured outputs,'' \textit{NIPS}, 2012.
\bibitem{rupprecht2017} C. Rupprecht et al., ``Learning in an uncertain world: Representing ambiguity through multiple hypotheses,'' \textit{ICCV}, 2017.
\bibitem{ijspeert2013} A. J. Ijspeert, J. Nakanishi, H. Hoffmann, P. Pastor, and S. Schaal, ``Dynamical movement primitives: Learning attractor models for motor behaviors,'' \textit{Neural Computation}, 2013.
\bibitem{paraschos2013} A. Paraschos, C. Daniel, J. Peters, and G. Neumann, ``Probabilistic movement primitives,'' \textit{NIPS}, 2013.
\bibitem{piegl1997} L. Piegl and W. Tiller, \textit{The NURBS Book}, 2nd ed., Springer, 1997.
\end{thebibliography}
}

\end{document}
